\section{Fuentes y nivel de lectura}
\label{ap:fuentes}

Casi todo el sustento bibliográfico de este informe proviene de una etapa de literatura previa, cuyos detalles están en \texttt{PROVENANCE.md}. Lo que sigue dice, para cada fuente, cuánto se leyó realmente, porque una afirmación apoyada en un resumen no tiene el mismo peso que una apoyada en el texto. Los niveles son: \textbf{texto parcial}, se extrajo el texto del PDF y se leyeron las ecuaciones y secciones indicadas (no el artículo entero); \textbf{resumen}, solo se verificaron el resumen y los metadatos, de modo que lo que se dice de ese artículo es lo que su resumen afirma y debería releerse antes de usarlo como sustento; \textbf{no leído}, no se accedió al texto (de pago o no descargado). Las citas textuales que aparecen en este informe son traducciones libres, hechas por quien escribió el texto, de los originales en inglés.

\begingroup\small
\renewcommand{\arraystretch}{1.2}
\begin{longtable}{@{}L{4.2cm}L{2.6cm}L{7.2cm}@{}}
\toprule
\textbf{Fuente} & \textbf{Nivel} & \textbf{Qué se toma}\\
\midrule
\endhead
León (2022) \cite{leon22} & texto parcial & Ecs.\ de fondo con $Q$ (9)--(12), (15), (57); espectro escalar (66); $N_f$; ausencia de sector tensorial\\
Piccirilli y León (2023) \cite{piccirilli23} & texto parcial & Ecs.\ (38), (40)--(42), (44); $N_f=100$; $r=16\epsilon_1$ sin derivación\\
Bengochea et al.\ (2023) \cite{bengochea23} & texto parcial & Acción (2.1), ecuaciones (2.5)--(2.15), tiempo unimodular (3.12), invariancia de gauge (5.10)\\
Padilla y Saltas (2015) \cite{padilla15} & texto parcial & Acción con multiplicador; equivalencia clásica con RG; acción de Henneaux--Teitelboim\\
Ellis et al.\ (2011, 2014) \cite{ellis11,ellis14} & resumen (+ búsqueda negativa de tensores) & Ecuaciones sin traza; no resuelven el valor de $\Lambda$; inflación ``como de costumbre''\\
Carballo-Rubio et al.\ (2022) \cite{carballo22} & texto parcial & Teoría lineal: WTDiff frente a Fierz--Pauli\\
Basak et al.\ (2016) \cite{basak16} & texto parcial & Resumen y conclusiones: tensores no afectados por el vínculo\\
Gao et al.\ (2014) \cite{gao14} & texto parcial & Restricción del gauge escalar; solo sector escalar\\
Cho y Singh (2015) \cite{cho15} & texto parcial & Acción tensorial independiente de $\xi$; $P_T$ estándar\\
Fabris et al.\ (2022) \cite{fabris22} & texto parcial & Ondas gravitacionales en UG no conservativa; Ec.~(37)\\
Chakraborty et al.\ (2025) \cite{chakraborty25} & texto parcial & Suma de polarizaciones\\
Barvinsky y Kolganov (2019) \cite{barvinsky19} & texto parcial & Acción tensorial en la UG generalizada, Ec.~(4.8)\\
Alvarenga et al.; Linares Cedeño y Nucamendi \cite{alvarenga23,linares23} & resumen & Diferencias en el sector escalar\\
Josset et al.; Perez y Sudarsky; Landau et al.; Cataldo \cite{josset17,perez19,landau23,cataldo26} & resumen / metadatos & Motivación física de $Q$; confrontación con datos; condición $\dot Q<0$\\
Henneaux y Teitelboim; Unruh; Weinberg \cite{henneaux89,unruh89,weinberg89} & no leídos & Solo por referencia; la forma de la acción de \cite{henneaux89} se verificó a través de \cite{padilla15}\\
Anderson y Finkelstein; Ng y van Dam \cite{anderson71,ng91} & no leídos & Solo referencias históricas\\
Martin, Ringeval y Vennin (2014) \cite{martin14} & texto parcial & Ecs.\ (2.17), (2.19)--(2.25): $P_h$, coeficientes de slow-roll, $C$\\
Baumann (2022) y TASI \cite{baumann22} & texto parcial & Ecs.\ (8.127)--(8.128); TASI arXiv:0907.5424v2, ecs. 196 y 198: solución de de Sitter a tiempo finito\\
Planck 2018 X \cite{planck20} & texto parcial & Cotas de $r$; valor de $A_s$; el enunciado sobre el potencial cuadrático\\
BICEP/Keck (2021) \cite{bicep21} & resumen & Cota $r<0.036$\\
Mukhanov, Feldman y Brandenberger (1992) \cite{mfb92} & no leído & Referencia de perturbaciones (pendiente de contrastar los pasos estándar)\\
Garriga y Mukhanov (1999) \cite{garriga99} & resumen & Referencia para la velocidad del sonido (hipótesis no verificada)\\
\bottomrule
\end{longtable}
\endgroup

Vale la pena explicitar una limitación: la búsqueda bibliográfica se hizo con una herramienta de búsqueda web, con la API de arXiv y con las bibliografías de los artículos de León, Piccirilli y Bengochea, y puede haber omitido trabajos sin arXiv. Del total de unos treinta artículos relevantes, solo se leyó el texto de unos diez.
